% =============================================================================
% KAPITEL 2: THEORETISCHE GRUNDLAGEN
% =============================================================================

\chapter{Theoretische Grundlagen}\label{chap:grundlagen}

Dieses Kapitel legt die theoretischen Fundamente für die Vorhersage von
Bundesliga-Spielergebnissen mittels Machine Learning. Abschnitt~\ref{sec:ml-grundlagen}
führt in Machine Learning als überwachtes Klassifikationsproblem ein und
beleuchtet die Besonderheiten des Sportkontexts. Abschnitt~\ref{sec:modelle}
stellt die fünf eingesetzten Algorithmen vor. Abschnitt~\ref{sec:vorhersagbarkeit}
diskutiert die grundlegenden Grenzen der Vorhersagbarkeit multivariater
Ereignisse sowie die theoretischen Konzepte für den Umgang mit ihnen.
Abschnitt~\ref{sec:dimensionsreduktion} behandelt Verfahren der
Dimensionsreduktion und Feature-Selektion. Abschnitt~\ref{sec:expected-goals}
schließt mit dem Konzept der Expected Goals ab.

% =============================================================================
\section{Machine Learning als überwachtes Klassifikationsproblem}
\label{sec:ml-grundlagen}
% =============================================================================

\textbf{Machine Learning} (ML) bezeichnet einen Teilbereich der Künstlichen
Intelligenz, bei dem Systeme aus Daten lernen, ohne für jede Entscheidung
explizit programmiert zu werden \cite{JamesEtAl2021}. Die Literatur
unterscheidet drei Lernparadigmen: überwachtes, unüberwachtes und
bestärkendes Lernen \cite{JamesEtAl2021,BunkerThabtah2019}. Die vorliegende
Arbeit befasst sich ausschließlich mit \textbf{überwachtem Lernen}: Der
Algorithmus erhält eine Menge von Trainingsbeispielen $\{(\mathbf{x}_i,
y_i)\}_{i=1}^{n}$ und lernt eine Funktion $f: \mathbf{x} \mapsto y$,
die für neue, ungesehene Eingaben $\mathbf{x}$ möglichst korrekte
Ausgaben $\hat{y}$ produziert.

\subsection{Multiclass-Klassifikation}

Innerhalb des überwachten Lernens ist die Spielvorhersage ein
\textbf{Multiclass-Klassifikationsproblem}: Die Zielvariable $y$ nimmt
eine von $K = 3$ diskreten Klassen an -- Heimsieg (\texttt{H}),
Unentschieden (\texttt{D}) oder Auswärtssieg (\texttt{A}). Im Gegensatz
zur binären Klassifikation muss ein Multiclass-Klassifikator gleichzeitig
zwischen mehreren Klassen unterscheiden. Dies erfordert entweder eine
inhärente Multiclass-Unterstützung des Algorithmus (wie bei der Softmax-
Logistischen Regression) oder eine Zerlegungsstrategie wie
\textit{One-vs-Rest} \cite{JamesEtAl2021,HastieEtAl2009}.

Ein besonderes Problem im Fußballkontext ist die \textbf{Klassen-Imbalance}:
Heimsiege machen in der Bundesliga etwa 45\,\% aller Ergebnisse aus,
Auswärtssiege ca.\ 30\,\% und Unentschieden nur rund 25\,\% \cite{HeGarcia2009}.
Unausgewogene Klassenverteilungen veranlassen viele Algorithmen dazu, die
Mehrheitsklasse zu bevorzugen und die Minderheitsklasse systematisch zu
ignorieren \cite{JapkowiczStephen2002}. Dieser Effekt lässt sich durch
drei Strategien abmildern:

\begin{itemize}
    \item \textbf{Class Weights:} Den Klassen werden beim Training umgekehrt
          proportionale Gewichte zugewiesen. Ein seltenes Unentschieden wird
          damit stärker bestraft, wenn es falsch klassifiziert wird
          \cite{JamesEtAl2021}.
    \item \textbf{Sample Weights:} Statt klassenspezifischer Gewichte wird
          jedem Trainingsbeispiel individuell ein Gewicht zugewiesen, das
          invers zur Häufigkeit seiner Klasse ist \cite{ScikitLearn2024}.
          Dies ist insbesondere bei Algorithmen nötig, die keine direkte
          \texttt{class\_weight}-Option unterstützen.
    \item \textbf{Oversampling:} Die Minderheitsklassen werden durch
          Wiederholung oder synthetische Generierung (z.\,B.\ SMOTE) auf
          die Größe der Mehrheitsklasse aufgefüllt \cite{ChawlaEtAl2002}.
          Dabei darf das Oversampling ausschließlich auf dem Trainingsdatensatz
          stattfinden, um eine Kontamination des Testsets zu vermeiden.
\end{itemize}

\subsection{Bias-Varianz-Dilemma, Overfitting und Underfitting}

Ein zentrales theoretisches Konzept des maschinellen Lernens ist das
\textbf{Bias-Varianz-Dilemma} \cite{GemanEtAl1992}. Der erwartete
quadratische Fehler eines Modells lässt sich zerlegen in:
\begin{equation}
    \mathbb{E}[(y - \hat{f}(\mathbf{x}))^2]
    = \underbrace{\text{Bias}^2[\hat{f}]}_{\text{Underfitting}}
    + \underbrace{\text{Var}[\hat{f}]}_{\text{Overfitting}}
    + \underbrace{\sigma^2}_{\text{irreduzibel}}
\end{equation}
\textbf{Bias} beschreibt systematische Fehler, die entstehen, wenn das
Modell zu simpel ist, um die wahren Zusammenhänge in den Daten zu
erfassen (\textit{Underfitting}). \textbf{Varianz} beschreibt die
Überempfindlichkeit eines Modells gegenüber den spezifischen
Trainingsdaten: Das Modell lernt deren Rauschen auswendig, anstatt
verallgemeinerbare Muster zu extrahieren (\textit{Overfitting}). Der
Term $\sigma^2$ ist der \textbf{irreduzible Fehler} -- derjenige Anteil
der Zielvariablen-Varianz, der auf nicht beobachtbare Einflüsse zurückgeht
und durch kein Modell eliminiert werden kann \cite{JamesEtAl2021}.

Overfitting ist besonders bei kleinen Datensätzen gefährlich. Mit nur
rund 306 Spielen pro Bundesliga-Saison und insgesamt etwa 2\,900
Beobachtungen ist die Gefahr, dass komplexe Modelle die Trainingsdaten
auswendig lernen, anstatt echte Muster zu erkennen, nicht trivial
\cite{Sumpter2016}. Techniken zur Kontrolle des Overfittings umfassen
\textbf{Regularisierung} (L1/L2-Strafterme), \textbf{Early Stopping}
(Abbruch des Trainings, sobald die Validierungsperformance nicht mehr
steigt) \cite{GoodfellowEtAl2016} sowie \textbf{Cross-Validation} zur
robusten Schätzung der Generalisierungsleistung.

\subsection{Modellselektion und Hyperparameteroptimierung}

Alle ML-Algorithmen besitzen \textbf{Hyperparameter} -- Steuergrößen,
die nicht aus den Daten gelernt, sondern vor dem Training vom Anwender
gesetzt werden (z.\,B.\ Baumtiefe, Lernrate, Regularisierungsstärke).
Die Suche nach guten Hyperparameter-Kombinationen wird als
\textbf{Hyperparameteroptimierung} bezeichnet \cite{JamesEtAl2021}.
Da eine vollständige Rastersuche (\textit{Grid Search}) bei großen
Suchräumen exponentiell teuer wird, bietet die \textbf{Randomized Search}
eine effiziente Alternative: Sie sampelt zufällig eine festgelegte Anzahl
von Kombinationen aus dem Suchraum und identifiziert so häufig annähernd
optimale Konfigurationen bei deutlich geringerem Rechenaufwand
\cite{JamesEtAl2021}.

Um zu verhindern, dass die Hyperparameter auf dem Testset überoptimiert
werden (\textit{Overfitting auf Hyperparameter-Ebene}), erfolgt die
Selektion stets über \textbf{Kreuzvalidierung}. Bei der
\textbf{Stratified K-Fold Cross-Validation} wird der Trainingsdatensatz
in $K$ gleichgroße Faltungen aufgeteilt, wobei die Klassenverteilung
in jeder Faltung der des Gesamtdatensatzes entspricht (\textit{Stratifizierung}).
Für jede der $K$ Runden dient eine andere Faltung als Validierungsmenge,
die restlichen $K-1$ als Training. Das Ergebnis ist eine stabile
Performance-Schätzung, die robust gegenüber Zufallsschwankungen in
einzelnen Datensplits ist \cite{JamesEtAl2021}.

\subsection{Besonderheiten von Sportdaten im ML-Kontext}

Sportdaten im Allgemeinen und Fußballdaten im Besonderen weisen
charakteristische Eigenschaften auf, die bei der Modellierung berücksichtigt
werden müssen \cite{BunkerThabtah2019,Sumpter2016}:

\begin{itemize}
    \item \textbf{Temporale Abhängigkeiten:} Die Leistung eines Teams in
          Spielwoche $t$ ist nicht unabhängig von der Leistung in Woche
          $t-1$. Features müssen diese zeitliche Struktur abbilden, ohne
          dabei zukünftige Information einzubeziehen (\textit{Data Leakage}).
    \item \textbf{Klassen-Imbalance:} Wie oben beschrieben sind die drei
          Ergebnisklassen ungleich häufig vertreten \cite{HeGarcia2009}.
    \item \textbf{Hohe Varianz und Rauschen:} Fußball ist ein
          Low-Scoring-Sport mit starkem Zufallseinfluss. Ein zufälliger
          Abpraller kann ein Spiel entscheiden, das eine Mannschaft
          dominiert hat \cite{DobsonGoddard2011}.
    \item \textbf{Multikollinearität:} Viele Spielstatistiken (z.\,B.
          Schüsse und Schüsse aufs Tor, xG und Tore) korrelieren stark
          miteinander, was Schätzungen in linearen Modellen destabilisiert
          \cite{JamesEtAl2021}.
    \item \textbf{Datenqualität und -integration:} Spielstatistiken
          stammen oft aus unterschiedlichen Quellen mit inkonsistenten
          Teamnamen, fehlenden Werten oder unterschiedlichen
          Erfassungsstandards. Eine sorgfältige Normalisierung,
          Qualitätsprüfung und Zusammenführung der Quellen ist
          Grundvoraussetzung für valide Modelle \cite{BunkerThabtah2019}.
\end{itemize}

% =============================================================================
\section{Eingesetzte Machine-Learning-Modelle}
\label{sec:modelle}
% =============================================================================

Tabelle~\ref{tab:model-uebersicht} gibt einen Überblick über die fünf
in dieser Arbeit evaluierten Algorithmen. Die Auswahl deckt ein breites
Spektrum von interpretablen linearen Modellen über Ensemble-Methoden bis
hin zu neuronalen Netzen ab.

\begin{table}[h]
    \centering
    \caption{Übersicht der eingesetzten Machine-Learning-Modelle}
    \label{tab:model-uebersicht}
    \begin{tabular}{l l l l}
        \toprule
        \textbf{Modell} & \textbf{Typ} & \textbf{Komplexität} & \textbf{Interpretierbarkeit} \\
        \midrule
        Logistische Regression & Linear          & Niedrig   & Hoch \\
        Random Forest          & Bagging-Ensemble & Mittel    & Mittel \\
        LightGBM               & Boosting-Ensemble & Hoch     & Mittel \\
        XGBoost                & Boosting-Ensemble & Hoch     & Mittel \\
        Multilayer Perceptron  & Neuronales Netz  & Sehr hoch & Niedrig \\
        \bottomrule
    \end{tabular}
\end{table}

\subsection{Logistische Regression}
\label{sec:logreg}

Die \textbf{Logistische Regression} ist das klassische lineare Modell
für Klassifikationsaufgaben und dient in dieser Arbeit als
\textbf{Baseline} \cite{JamesEtAl2021,HastieEtAl2009}. Für $K = 3$
Klassen modelliert sie die Klassenwahrscheinlichkeiten über die
Softmax-Funktion:
\begin{equation}
    P(y=k \mid \mathbf{x}) = \frac{
        \exp\!\left(\beta_{k0} + \sum_{j=1}^{p} \beta_{kj} x_j\right)
    }{
        \sum_{l=1}^{K} \exp\!\left(\beta_{l0} + \sum_{j=1}^{p} \beta_{lj} x_j\right)
    }
\end{equation}
Die Koeffizienten $\beta_{kj}$ werden durch Maximierung der
Log-Likelihood geschätzt. \textbf{L2-Regularisierung} (Ridge) mit dem
inversen Strafterm $C$ kontrolliert Overfitting und verhindert zu große
Koeffizientenwerte \cite{HastieEtAl2009}. Die logistische Regression
bietet hohe Interpretierbarkeit -- die Koeffizientenvorzeichen geben
direkt an, welche Features einen Heimsieg begünstigen oder hemmen --
bei gleichzeitig geringem Rechenaufwand \cite{BunkerThabtah2019}.
Nachteilig ist ihre Unfähigkeit, nicht-lineare Zusammenhänge und
Feature-Interaktionen ohne manuelle Erweiterungen zu erfassen.

\subsection{Random Forest}
\label{sec:random-forest}

Der \textbf{Random Forest} ist ein Ensemble-Verfahren, das nach dem
Prinzip des \textbf{Bootstrap Aggregating (Bagging)} eine Vielzahl von
Entscheidungsbäumen kombiniert \cite{Breiman2001}. Jeder Baum wird auf
einer Zufallsstichprobe der Trainingsdaten (mit Zurücklegen) trainiert.
Zusätzlich wird bei jedem Knoten-Split nur eine zufällige Teilmenge
von $m \approx \sqrt{p}$ Features als Kandidaten zugelassen. Diese
doppelte Randomisierung dekorreliert die Bäume, was die Varianz des
Ensembles gegenüber einzelnen Bäumen erheblich reduziert, ohne den Bias
wesentlich zu erhöhen \cite{JamesEtAl2021}.

Ein wichtiger Nebeneffekt ist die \textbf{Feature Importance}: Für
jeden Baum und jeden Split kann gemessen werden, wie stark die jeweilige
Feature-Variable die Reinheit der entstehenden Knoten verbessert
(\textit{Gini Importance}). Über alle Bäume gemittelt ergibt dies eine
stabile Aussage darüber, welche Features für die Vorhersage am
einflussreichsten sind \cite{Breiman2001,LundbergLee2017}.

\subsection{Gradient Boosting: XGBoost und LightGBM}
\label{sec:gradient-boosting}

\textbf{Gradient Boosting} \cite{Friedman2001} ist ein sequenzielles
Ensemble-Verfahren: Jeder neue Baum $f_m$ wird darauf trainiert, die
Residuen (genauer: den negativen Gradienten der Verlustfunktion) seines
Vorgänger-Ensembles zu korrigieren. Das finale Ensemble lautet:
\begin{equation}
    F_M(\mathbf{x}) = \sum_{m=1}^{M} \eta \cdot f_m(\mathbf{x})
\end{equation}
wobei $\eta$ die \textit{Lernrate} ist, die bestimmt, wie stark jeder
neue Baum in das Ensemble einfließt. Eine kleine Lernrate erfordert
mehr Bäume ($M$ groß), erzeugt aber in der Regel robustere Modelle
\cite{JamesEtAl2021}.

\textbf{XGBoost} \cite{ChenGuestrin2016} erweitert das Gradient-Boosting-
Rahmenwerk um eine regularisierte Verlustfunktion, die sowohl L1- als
auch L2-Strafterme enthält, sowie um System-Optimierungen für
paralleles Training. Zusätzliche Hyperparameter wie \texttt{gamma}
(Mindestschwelle für einen Knoten-Split) und \texttt{min\_child\_weight}
(Mindestanzahl von Trainingsbeispielen pro Blatt) ermöglichen eine
feingranulare Kontrolle des Overfitting-Risikos.

\textbf{LightGBM} \cite{KeEtAl2017} verfolgt einen alternativen
Ansatz: Statt des üblichen \textit{level-wise} Baumwachstums (alle
Knoten einer Tiefe werden gleichzeitig gespalten) setzt LightGBM auf
\textit{leaf-wise} Wachstum -- der Knoten mit dem größten Gewinn wird
zuerst gespalten, unabhängig von seiner Tiefe. Hinzu kommen
\textit{Gradient-based One-Side Sampling} (GOSS), das Trainingsdaten
mit kleinem Gradienten -- also bereits gut gelernte Beispiele --
reduziert, und \textit{Exclusive Feature Bundling} (EFB), das
gegenseitig ausschließende Features zusammenfasst. Das Ergebnis ist
ein erheblich schnellerer und speichereffizienterer Algorithmus bei
vergleichbarer Qualität \cite{KeEtAl2017}.

Beide Gradient-Boosting-Implementierungen liefern ebenfalls
\textbf{Feature-Importance-Werte} und gehören zum State of the Art
für tabellarische Daten \cite{GrollLeyEffenberger2019,HubacekSourekZelezny2019}.

\subsection{Multilayer Perceptron}
\label{sec:mlp}

Das \textbf{Multilayer Perceptron} (MLP) ist ein feedforward neuronales
Netz, das aus einer Eingabeschicht, einer oder mehreren versteckten
Schichten (\textit{Hidden Layers}) und einer Ausgabeschicht besteht
\cite{GoodfellowEtAl2016}. Jede Schicht berechnet eine affine
Transformation der Eingabe, gefolgt von einer nicht-linearen
Aktivierungsfunktion:
\begin{equation}
    \mathbf{h}^{(l)} = \sigma\!\left(\mathbf{W}^{(l)} \mathbf{h}^{(l-1)} + \mathbf{b}^{(l)}\right)
\end{equation}
In den versteckten Schichten wird typischerweise \textbf{ReLU}
($\sigma(z) = \max(0, z)$) als Aktivierungsfunktion verwendet, da
sie das Problem verschwindender Gradienten bei tiefen Netzen abmildert.
Die Ausgabeschicht verwendet \textbf{Softmax}, um Klassenwahrscheinlichkeiten
zu erzeugen. Das Training erfolgt via \textbf{Backpropagation} und
stochastischem Gradientenabstieg, wobei die \textbf{Cross-Entropy-Loss}
minimiert wird \cite{RumelhartEtAl1986}.

Overfitting ist beim MLP besonders ausgeprägt, da die Zahl der
Parameter die Anzahl der Trainingsbeispiele leicht übersteigen kann.
Gegenmaßnahmen umfassen \textbf{L2-Gewichtsregularisierung},
\textbf{Dropout} (zufälliges Deaktivieren von Neuronen während des
Trainings \cite{SrivastavaEtAl2014}) sowie \textbf{Early Stopping}:
Das Training wird abgebrochen, sobald der Validierungsverlust über
mehrere Epochen hinweg nicht weiter sinkt \cite{GoodfellowEtAl2016}.
Der bis dahin beste Modellzustand wird gespeichert. Early Stopping
wirkt gleichzeitig als implizite Regularisierung und erspart die
manuelle Wahl der Epochenanzahl.

% =============================================================================
\section{Vorhersagbarkeit multivariater Ereignisse}
\label{sec:vorhersagbarkeit}
% =============================================================================

\subsection{Das multivariate Vorhersageproblem}

Das Ergebnis eines Fußballspiels hängt von einer theoretisch unbegrenzten
Anzahl von Einflussgrößen ab: Spielerform, Verletzungen, Tagesform,
Wetter, Schiedsrichterentscheidungen, taktische Anpassungen während des
Spiels, Platzierungsdruck, Reisemüdigkeit und viele weitere Faktoren
\cite{DobsonGoddard2011,Sumpter2016}. Formal lässt sich dies schreiben als:
\begin{equation}
    Y = f(X_1, X_2, \dots, X_p, X_{p+1}, \dots) + \varepsilon
\end{equation}
wobei $X_1, \dots, X_p$ die beobachtbaren Features sind und alle weiteren
Einflussgrößen im Störterm $\varepsilon$ gebündelt werden. Dieser enthält
den \textbf{irreduziblen Fehler} $\sigma^2$ -- denjenigen Varianzanteil,
der prinzipiell nicht durch ein Modell erklärt werden kann, weil die
zugrundeliegenden Faktoren nicht messbar oder modellierbar sind
\cite{GemanEtAl1992}. Selbst ein theoretisch perfektes Modell mit Zugang
zu allen beobachtbaren Prädiktoren würde diesen Fehler nicht beseitigen
können.

\subsection{Feature Engineering und zeitabhängige Formabbildung}

Da Rohdaten (Tabellenpositionen, Einzelspiel-Tore) kurzfristige
Schwankungen schlecht abbilden, sind aufbereitete \textbf{Form-Features}
zentral für die Vorhersagequalität. Zwei Konzepte sind dabei besonders
relevant:

\textbf{Rolling Windows} berechnen Durchschnitte oder Summen über ein
gleitendes Zeitfenster der letzten $n$ Spiele \cite{BunkerThabtah2019}.
Ein einfaches Rolling Window gewichtet alle Spiele innerhalb des Fensters
gleich. Dieser Ansatz ist intuitiv, reagiert jedoch sprunghaft, wenn ein
altes Spiel aus dem Fenster fällt.

\textbf{Exponentially Weighted Averages (EWA)} lösen dieses Problem,
indem neueren Spielen exponentiell höheres Gewicht zugewiesen wird
\cite{Brown1956,HyndmanAthanasopoulos2021}. Der EWA zum Zeitpunkt $t$
berechnet sich rekursiv als:
\begin{equation}
    S_t = \alpha \cdot x_t + (1 - \alpha) \cdot S_{t-1},
    \quad \alpha = \frac{2}{s + 1}
\end{equation}
wobei $s$ der Span-Parameter ist und $\alpha \in (0, 1)$ der
Glättungsfaktor. Ein großes $s$ (flache Gewichtung) ergibt einen
stabilen, aber trägen Durchschnitt; ein kleines $s$ (steile Gewichtung)
reagiert schnell auf aktuelle Form, ist aber anfälliger für Rauschen.
EWA gelten in der Zeitreihenanalyse als robuster als einfache Rolling
Windows, da sie ohne scharfe Fensterränder auskommen \cite{HyndmanAthanasopoulos2021}.

\subsection{Data Leakage}

\textbf{Data Leakage} bezeichnet das unbeabsichtigte Einfließen von
Information aus der Zukunft in die Modellierung \cite{KaufmanRossetPerlich2012}.
Im Sportkontext tritt es auf, wenn zum Beispiel der EWA-Wert eines Spiels
das Ergebnis dieses Spiels selbst einschließt, oder wenn saisonale
Statistiken zum Zeitpunkt eines Spiels noch nicht feststehen. Data Leakage
führt zu künstlich hohen Trainingsperformances und zu Modellen, die in
der Praxis deutlich schlechter abschneiden als in der Evaluation -- da
in der Praxis die \glqq zukunftige\grqq{} Information schlicht nicht
verfügbar ist \cite{KaufmanRossetPerlich2012}. Die konsequente Vermeidung
von Data Leakage ist eine der wichtigsten methodischen Anforderungen an
zeitreihenbasierte Vorhersagesysteme.

\subsection{Ausreißer und IQR-Bereinigung}

\textbf{Ausreißer} sind Beobachtungen, die stark vom typischen Muster
abweichen. Sie können durch Datenfehler, außergewöhnliche Spielverläufe
(z.\,B.\ 9:0-Ergebnis) oder seltene Ereignisse entstehen. Ausreißer
können Modelle verzerren, insbesondere lineare Modelle und Algorithmen,
die empfindlich auf die Skalierung der Eingaben reagieren
\cite{JamesEtAl2021}.

Eine verbreitete Methode zur Ausreißererkennung ist der
\textbf{Interquartilsabstand (IQR)}. Für ein gegebenes Feature gilt:
\begin{equation}
    \text{IQR} = Q_3 - Q_1
\end{equation}
Beobachtungen außerhalb des Intervalls $[Q_1 - c \cdot \text{IQR},\;
Q_3 + c \cdot \text{IQR}]$ werden als Ausreißer klassifiziert. Der
Multiplikator $c$ bestimmt die Strenge der Bereinigung: $c = 1{,}5$
ist der Standardwert für moderate Bereinigung, $c = 3{,}0$ für die
Entfernung nur extremer Ausreißer \cite{JamesEtAl2021}. Der Vorteil
des IQR gegenüber z-Score-basierten Methoden ist die Robustheit
gegenüber schiefen Verteilungen, da er kein Gauß'sches Grundmodell
voraussetzt.

\subsection{Korrelation und Pearson-Korrelationskoeffizient}

Der \textbf{Pearson-Korrelationskoeffizient} $r$ misst den linearen
Zusammenhang zwischen zwei Variablen \cite{JamesEtAl2021}:
\begin{equation}
    r_{XY} = \frac{\sum_{i=1}^{n}(x_i - \bar{x})(y_i - \bar{y})}
                  {\sqrt{\sum_{i=1}^{n}(x_i - \bar{x})^2 \cdot
                         \sum_{i=1}^{n}(y_i - \bar{y})^2}}
\end{equation}
$r$ liegt stets im Intervall $[-1, 1]$: $r = 1$ entspricht perfekter
positiver Korrelation, $r = -1$ perfekter negativer Korrelation,
$r = 0$ linearer Unabhängigkeit. Im Feature-Engineering-Kontext
wird der Pearson-Koeffizient genutzt, um die Stärke des Zusammenhangs
jedes Features mit der Zielvariablen zu beurteilen. Zusätzlich hilft
die Korrelationsmatrix aller Features, stark korrelierte Prädiktoren
zu identifizieren (\textbf{Multikollinearität}), die redundante
Information tragen und lineare Modelle destabilisieren können
\cite{JamesEtAl2021,HastieEtAl2009}.

\subsection{Standardisierung}

Viele Algorithmen -- insbesondere Logistische Regression, SVM und MLP
-- sind empfindlich auf unterschiedliche Skalen der Features: Ein Feature
mit Werten im Bereich $[0, 1000]$ dominiert numerisch gegenüber einem
Feature im Bereich $[0, 1]$, obwohl letzteres inhaltlich genauso wichtig
sein kann \cite{JamesEtAl2021}. \textbf{Standardisierung} (auch
Z-Score-Normierung) transformiert jedes Feature auf Mittelwert 0 und
Standardabweichung 1:
\begin{equation}
    \tilde{x}_j = \frac{x_j - \mu_j}{\sigma_j}
\end{equation}
Dabei müssen $\mu_j$ und $\sigma_j$ ausschließlich aus den
\textbf{Trainingsdaten} geschätzt werden. Die Anwendung auf die
Testdaten erfolgt mit diesen Werten, um eine Datenkontamination zu
vermeiden \cite{HastieEtAl2009}.

\subsection{Train-Test-Split}

Die grundlegende Strategie zur unverzerrten Schätzung der
Generalisierungsleistung ist die Aufteilung des Datensatzes in einen
\textbf{Trainingsdatensatz} (zum Modelltraining) und einen \textbf{Testdatensatz}
(zur abschließenden, einmaligen Evaluation) \cite{JamesEtAl2021}.
Um sicherzustellen, dass beide Teilmengen repräsentative Klassenverteilungen
aufweisen, wird eine \textbf{stratifizierte} Aufteilung verwendet: Der
Anteil von Heimsiegen, Unentschieden und Auswärtssiegen ist in Trainings-
und Testset identisch. Dies ist bei Klassen-Imbalance besonders wichtig,
da andernfalls die seltene Klasse im Testset unter- oder überrepräsentiert
sein könnte \cite{JamesEtAl2021}.

% =============================================================================
\section{Dimensionsreduktion und Feature-Selektion}
\label{sec:dimensionsreduktion}
% =============================================================================

\subsection{Der Fluch der Dimensionalität}

Mit steigender Anzahl von Features wächst der Merkmalsraum exponentiell.
\citet{Bellman1961} prägte dafür den Begriff des \textbf{Curse of
Dimensionality}: In hochdimensionalen Räumen werden Datenpunkte
zunehmend gleichweit voneinander entfernt, was Distanzmetriken und
Ähnlichkeitsmaße ihre Aussagekraft verlieren lässt \cite{HastieEtAl2009}.
Für die Fußballvorhersage bedeutet das: Mehr Features sind nicht
automatisch besser. Redundante oder schwach prädiktive Features
erhöhen die Modellkomplexität, fördern Overfitting und erschweren die
Interpretation \cite{JamesEtAl2021}.

\subsection{Feature-Selektion}

\textbf{Feature-Selektion} wählt eine Teilmenge der ursprünglichen
Features aus und behält dabei die Interpretierbarkeit der einzelnen
Variablen \cite{GuyonElisseeff2003}. Drei Hauptkategorien lassen sich
unterscheiden:

\begin{itemize}
    \item \textbf{Filter-Methoden} bewerten Features modellunabhängig
          anhand statistischer Maße wie Pearson-Korrelation, Mutual
          Information oder dem ANOVA-F-Test \cite{Fisher1935,CoverThomas2006}.
          Sie sind schnell und transparent, berücksichtigen aber keine
          Feature-Interaktionen.
    \item \textbf{Wrapper-Methoden} wählen Features iterativ aus, indem
          sie die tatsächliche Modellperformance als Auswahlkriterium
          verwenden (z.\,B.\ Recursive Feature Elimination) \cite{KohaviJohn1997,GuyonEtAl2002}.
          Sie sind genauer, aber rechenintensiv.
    \item \textbf{Embedded-Methoden} integrieren die Selektion direkt in
          das Modelltraining, z.\,B.\ über L1-Regularisierung (\textit{Lasso}),
          die Koeffizienten unwichtiger Features auf null setzt
          \cite{Tibshirani1996}, oder über tree-basierte Feature-Importance-
          Werte \cite{Breiman2001,LundbergLee2017}.
\end{itemize}

\subsection{Principal Component Analysis (PCA)}

Die \textbf{Principal Component Analysis} (PCA) ist ein klassisches
lineares Verfahren zur Dimensionsreduktion \cite{Pearson1901,Hotelling1933}.
Im Gegensatz zur Feature-Selektion transformiert PCA die ursprünglichen
Features in einen neuen Merkmalsraum aus orthogonalen
\textbf{Hauptkomponenten}, die die Varianz der Daten maximal erklären
\cite{Jolliffe2002}. Die $k$-te Hauptkomponente $\mathbf{w}_k$ ist
die Richtung im Merkmalsraum, entlang derer die projizierte Varianz
maximal ist und die orthogonal zu allen vorherigen Komponenten steht:
\begin{equation}
    \mathbf{w}_k = \underset{\|\mathbf{w}\|=1,\; \mathbf{w} \perp \mathbf{w}_{1..k-1}}{\arg\max}\;
    \operatorname{Var}(\mathbf{X}\mathbf{w}), \qquad
    \boldsymbol{\Sigma}\, \mathbf{w}_k = \lambda_k\, \mathbf{w}_k
\end{equation}
In der Praxis wählt man die Anzahl der Komponenten $K$ so, dass ein
Mindestanteil der Gesamtvarianz (typischerweise $\geq 95\,\%$) erhalten
bleibt. PCA ist besonders wirksam bei stark korrelierten Features
(\textit{Multikollinearität}), da die ersten Hauptkomponenten gerade
die gemeinsame Varianz mehrerer korrelierter Originalfeatures bündeln.

Der Hauptnachteil der PCA ist der \textbf{Verlust der Interpretierbarkeit}:
Die resultierenden Komponenten sind Linearkombinationen aller ursprünglichen
Features und haben keine direkte inhaltliche Bedeutung mehr. Zudem
ist PCA sensitiv gegenüber Ausreißern und erfordert eine vorherige
Standardisierung \cite{Bishop2006}.

% =============================================================================
\section{Expected Goals -- Konzept und verwandte Metriken}
\label{sec:expected-goals}
% =============================================================================

\subsection{Das Konzept der Expected Goals}

\textbf{Expected Goals} (xG) ist eine fortgeschrittene Fußballstatistik,
die jeder Torschusssituation eine Wahrscheinlichkeit zwischen 0 und 1
zuweist -- die Wahrscheinlichkeit, dass aus dieser Situation ein Tor
erzielt wird \cite{Spearman2018,AnzerBauer2021}. Formal liegt ein
\textbf{binäres Klassifikationsproblem} vor:
\begin{equation}
    \text{xG}_i = P(\text{Tor}_i \mid \mathbf{x}_i) = f(\mathbf{x}_i),
    \qquad
    \text{xG}_{\text{Team}} = \sum_{i=1}^{n} \text{xG}_i
\end{equation}
Dabei ist $\mathbf{x}_i$ ein Merkmalsvektor, der die relevanten
Kontextfaktoren der Schusssituation $i$ beschreibt. Der Wert
$\text{xG}_{\text{Team}}$ ist die Summe aller Einzelchancen und
entspricht den \emph{erwarteten} Toren eines Teams in einem Spiel.

Die Entwicklung moderner xG-Modelle verlief in drei Phasen
\cite{Wheatcroft2022}: In der \textbf{Prä-xG-Ära} (vor 2010) wurden
ausschließlich Rohdaten wie Schussanzahl genutzt. Die \textbf{Entstehungsphase}
(2010--2015) führte erste logistische Regressionsmodelle auf Basis von
Schussdistanz und -winkel ein. Seit 2015 dominieren \textbf{moderne
xG-Modelle} auf Basis von Gradient Boosting, die AUC-Werte von über
0,85 erreichen \cite{AnzerBauer2021}.

\subsection{Einflussgrößen moderner xG-Modelle}

Die wesentlichen Feature-Kategorien moderner xG-Modelle umfassen
\cite{AnzerBauer2021,BauerEpstein2023,PalmerEtAl2021}:

\begin{itemize}
    \item \textbf{Geometrische Faktoren:} Schussdistanz zum Tor,
          Schusswinkel, Position im Strafraum. Ein Schuss aus dem
          Elfmeterpunkt hat einen xG-Wert von ca.\ 0,76
          \cite{PalmerEtAl2021}.
    \item \textbf{Körperteil:} Schüsse mit dem starken Fuß
          ($\sim$0,25 xG) erzielen im Durchschnitt höhere Werte als
          Kopfbälle ($\sim$0,15 xG) oder der schwache Fuß \cite{PalmerEtAl2021}.
    \item \textbf{Vorlagentyp:} Flanke, direkter Pass, Abpraller oder
          Elfmeter -- der Entstehungsweg beeinflusst die Abschlussqualität.
    \item \textbf{Druck- und Kontextfaktoren:} Anzahl verteidigender
          Spieler, Torwartposition, Spielstand.
\end{itemize}

\subsection{Verwandte Metriken: xGA und xPTS}

Neben xG sind zwei abgeleitete Metriken für die Teamleistungsanalyse
und Spielvorhersage relevant:

\textbf{Expected Goals Against (xGA)} ist das xG-Äquivalent auf
defensiver Seite: Es quantifiziert, wie viele Tore ein Team hätte
erwarten können zu kassieren, basierend auf den Schusssituationen,
die es zugelassen hat. xGA misst defensive Stabilität robuster als
die tatsächlich kassierten Tore, da es überragende Torwartleistungen
herausrechnet und den zufälligen Effekt von Wecheln der Effektivität
normalisiert \cite{VanEetveldeEtAl2024}.

\textbf{Expected Points (xPTS)} kombiniert xG und xGA auf Spielebene:
Aus den erwarteten Torverteilungen beider Teams wird simuliert, welchen
Punkteanteil jedes Team langfristig erwarten kann. xPTS ist ein
stabiler Indikator für die \glqq verdiente\grqq{} Ligaposition eines
Teams und identifiziert Over- und Underperformance zuverlässig
\cite{Wheatcroft2022}.

\subsection{Warum xG traditionelle Statistiken übertr\textit{i}fft}

Der zentrale Vorteil von xG gegenüber tatsächlichen Torzahlen liegt
in der \textbf{größeren Stichprobenbasis}: Während in einem Bundesliga-Spiel
durchschnittlich nur 2,86 Tore fallen, werden typischerweise 10 bis 20
Torchancen erzeugt. Jede Chance liefert eine xG-Schätzung, sodass die
Summe $\text{xG}_{\text{Team}}$ auf einer deutlich größeren Datenmenge
basiert als die bloße Toranzahl \cite{AnzerBauer2021}. Dies macht xG
stabiler gegenüber kurzfristigen Zufallsschwankungen und zu einem
besseren Prädiktor für nachhaltige Teamleistung
\cite{WheatcroftVanEetvelde2020,LiuEtAl2023}.

Ein klassisches Beispiel ist die Leicester-City-Meisterschaft 2015/16:
Traditionelle Statistiken spiegelten den Titelgewinn wider, doch xG
zeigte, dass Leicester deutlich über seinen erwarteten Werten performte.
In den Folgejahren pendelte sich die Teamleistung auf dem xG-entsprechenden
Niveau ein -- die Meisterschaft war teilweise ein statistischer Ausreißer
\cite{Wheatcroft2022}.

\subsection{Limitationen von xG}

Trotz seiner Stärken unterliegt xG mehreren \textbf{Limitationen}
\cite{Wheatcroft2022,BauerEpstein2023}: Erstens existiert kein einheitlicher
Goldstandard -- verschiedene Anbieter berechnen xG unterschiedlich,
was Vergleiche zwischen Datenquellen erschwert. Zweitens fehlen in
öffentlich verfügbaren Datensätzen häufig detaillierte Kontextfaktoren
wie Torwartposition oder Drucksituation. Drittens ist xG
\textbf{ligaspezifisch}: Ein auf Premier-League-Daten trainiertes
xG-Modell kann aufgrund unterschiedlicher Spielstile (die Bundesliga
weist traditionell eine höhere Torrate auf) für die Bundesliga
suboptimale Schätzungen liefern. Diese Einschränkungen unterstreichen
die Bedeutung ligaspezifisch kalibrierter xG-Daten.

% =============================================================================
\section{Evaluationsmetriken für Klassifikationsmodelle}
\label{sec:evaluationsmetriken}
% =============================================================================

Die Güte eines Klassifikationsmodells lässt sich anhand mehrerer
komplementärer Metriken beurteilen. Im Kontext der Fußballvorhersage
mit drei ungleich häufigen Klassen sind folgende Kennzahlen relevant.

\subsection{Grundlegende Metriken: Precision, Recall und F1-Score}

Ausgehend von der \textbf{Konfusionsmatrix} -- einer $K \times K$-Matrix,
die für jede Kombination aus tatsächlicher und vorhergesagter Klasse
die Anzahl der Beobachtungen zeigt -- lassen sich pro Klasse $k$ drei
Metriken ableiten \cite{Powers2020}:

\begin{itemize}
    \item \textbf{Precision$_k$:} Der Anteil der korrekt als Klasse $k$
          vorhergesagten Beobachtungen an allen als $k$ vorhergesagten:
          $\text{Prec}_k = \frac{TP_k}{TP_k + FP_k}$.
    \item \textbf{Recall$_k$:} Der Anteil der korrekt als Klasse $k$
          erkannten Beobachtungen an allen tatsächlich zur Klasse $k$
          gehörenden: $\text{Rec}_k = \frac{TP_k}{TP_k + FN_k}$.
    \item \textbf{F1-Score$_k$:} Das harmonische Mittel aus Precision
          und Recall: $F1_k = \frac{2 \cdot \text{Prec}_k \cdot \text{Rec}_k}{\text{Prec}_k + \text{Rec}_k}$.
\end{itemize}

\subsection{Macro-F1, Accuracy und Kalibrierung}

\textbf{Accuracy} ist der Anteil korrekt klassifizierter Beobachtungen
an allen Beobachtungen. Bei unausgewogenen Klassen ist sie irreführend:
Ein Modell, das immer \glqq Heimsieg\grqq{} vorhersagt, erzielt bereits
ca.\ 45\,\% Accuracy, ohne echtes Wissen über das Spiel zu haben
\cite{Powers2020}.

Der \textbf{Macro-F1-Score} ist das arithmetische Mittel der
klassenspezifischen F1-Scores:
\begin{equation}
    F1_{\text{macro}} = \frac{1}{K} \sum_{k=1}^{K} F1_k
\end{equation}
Er gewichtet alle Klassen gleich, unabhängig von ihrer Häufigkeit.
Ein Modell, das Unentschieden vollständig ignoriert, erhält einen
$F1_D = 0$ und damit einen deutlich niedrigeren Macro-F1, selbst wenn
Accuracy und Heimsieg-F1 hoch sind \cite{Powers2020,HeGarcia2009}.

\textbf{Kalibrierung} beschreibt, ob die vorhergesagten
Klassenwahrscheinlichkeiten mit den empirischen Häufigkeiten
übereinstimmen: Ein gut kalibriertes Modell, das einem Spiel eine
70\,\%-Heimsieg-Wahrscheinlichkeit zuweist, sollte in 70\,\% der
Fälle, in denen es diese Wahrscheinlichkeit vorhersagt, tatsächlich
ein Heimsieg eingetreten sein \cite{JamesEtAl2021}.

\subsection{Feature Importance und Modellstabilität}

\textbf{Feature Importance} beschreibt, wie stark einzelne Features
zur Modellentscheidung beitragen. Für baumbasierte Modelle wird sie
über den mittleren Gini-Gewinn über alle Splits und Bäume berechnet
\cite{Breiman2001}. Eine hohe Feature Importance eines xG-Features
würde empirisch belegen, dass xG-Informationen für die Vorhersage
entscheidend sind \cite{LundbergLee2017}.

\textbf{Modellstabilität} beschreibt die Konsistenz der
Modellperformance über verschiedene Datensplits. Stark schwankende
F1-Werte über die Faltungen der Kreuzvalidierung deuten auf
Overfitting oder eine geringe Datenbasis hin und sind ein Warnsignal
für mangelnde Generalisierbarkeit \cite{JamesEtAl2021}.

% =============================================================================
\section{Zusammenfassung}
\label{sec:zusammenfassung-grundlagen}
% =============================================================================

Dieses Kapitel hat die theoretischen Grundlagen für die Vorhersage von
Bundesliga-Spielergebnissen dargelegt:

\begin{itemize}
    \item \textbf{Machine Learning als überwachtes Klassifikationsproblem}
          mit drei Klassen und Klassen-Imbalance erfordert spezifische
          Techniken: class weights, sample weights oder oversampling sowie
          robuste Metriken wie Macro-F1 statt reiner Accuracy.
    \item \textbf{Overfitting und Underfitting} sind über das
          Bias-Varianz-Dilemma theoretisch fundiert; Early Stopping,
          Regularisierung und Kreuzvalidierung sind die zentralen
          Gegenmaßnahmen.
    \item \textbf{Fünf Algorithmen} -- Logistische Regression, Random
          Forest, XGBoost, LightGBM und MLP -- decken ein breites
          Spektrum von Modellkomplexität und Interpretierbarkeit ab.
    \item \textbf{Vorhersagbarkeit multivariater Ereignisse} ist durch
          irreduziblen Fehler, emergente Nicht-Linearitäten und fehlende
          Einflussvariablen fundamental begrenzt. EWA und Rolling Windows
          ermöglichen eine zeitvalide Formabbildung; Data Leakage muss
          konsequent vermieden werden.
    \item \textbf{Dimensionsreduktion} durch PCA und Feature-Selektion
          bekämpft den Curse of Dimensionality und reduziert Multikollinearität.
    \item \textbf{Expected Goals} (xG), xGA und xPTS sind robustere
          Leistungsindikatoren als reine Torstatistiken und bilden die
          Grundlage für das erweiterte Feature-Set dieser Arbeit.
\end{itemize}

Aufbauend auf diesen Grundlagen gibt Kapitel~\ref{chap:relatedwork}
einen Überblick über den aktuellen Stand der Forschung.